% ============================================================
\section{在 GUI 中显示地图}
\subsection*{Displaying Maps in the GUI}

分析完成后，工具将设计数据与分析结果（\texttt{*.result}）保存在 RedHawk 工作目录下的 \texttt{in-design.redhawk/design\_name.result} 目录中。必须先运行 \cmd{open\_rail\_result} 加载分析结果，然后才能在 Fusion Compiler GUI 中显示地图。

表~\ref{tab:map-types} 列出 RedHawk Fusion 分析流程中支持的地图类型。

\begin{table}[htbp]
\centering
\caption{RedHawk Fusion 流程中支持的分析地图（原书 Table 70）}
\label{tab:map-types}
\small
\begin{tabularx}{\textwidth}{@{}l X@{}}
\toprule
\textbf{地图类型} & \textbf{说明} \\
\midrule
Rail voltage drop map &
显示压降地图：叠加在物理电源 net 上的彩色编码压降值。静态分析显示平均压降；动态分析显示峰值压降、平均压降或峰值电压上升。
\begin{noteBox}
即使 Fusion Compiler 库文件中不存在某些层，工具仍在压降地图上显示 LEF 或 DEF 文件中的全部层。
\end{noteBox} \\
Rail parasitics map &
压降分析完成后，显示 block 中给定电源 net 的寄生电阻；显示该 net 任意给定形状的电阻。 \\
Rail power map &
压降分析完成后，基于 RedHawk 功耗计算结果显示基于实例的功耗地图或功耗密度地图。基于实例的功耗地图：显示 block 中单元实例的功耗值。功耗密度地图：显示 block 中区域的功耗密度值。 \\
Rail current map &
压降分析完成后显示 block 的电流分布。 \\
Rail minimal path resistance map &
最小路径电阻分析完成后，显示所选 net 的最小路径电阻值。 \\
Rail electromigration map &
静态或动态电迁移分析完成后，显示 block 的电流违例。 \\
Rail instance peak current map &
显示 block 中不同实例或区域的峰值电流值。 \\
Rail instance effective voltage drop map &
压降分析完成后，显示 block 中给定实例的有效电压值。 \\
Rail instance effective resistance map &
有效电阻分析完成后，显示 block 中给定实例的有效电阻值。 \\
Rail instance switch cell &
动态分析完成且开关单元模型输入文件可用时，显示 block 中开关单元的电流、有效电压与电压差地图。若无开关模型输入文件，工具将这些开关单元视为黑盒且不显示轨地图。基于实例的电流地图：显示电流如何流经所选开关单元。基于实例的有效电压地图：显示所选开关单元处的有效电压。基于实例的电压差地图：显示所选开关单元两端的电压差。 \\
\bottomrule
\end{tabularx}
\end{table}

要在 Fusion Compiler GUI 中显示地图：
\begin{enumerate}
  \item 要通过 GUI 加载轨分析结果，选择 Task $>$ Rail Analysis $>$ 2D Rail Analysis。出现 Block Level 2D Rail Analysis 窗口后，单击 \texttt{open\_rail\_result} 链接。

或者运行 \cmd{open\_rail\_result}，加载分析完成后保存在 RedHawk 工作目录下 \texttt{design\_name} 目录中的设计数据与分析结果（\texttt{*.result}）。
  \item 在 GUI 布局窗口中，选择 View $>$ Map 并选择要显示的地图。

在地图中，问题区域以不同颜色高亮。将指针移到实例上可在 InfoTip 中查看更多信息。默认工具显示 block 中全部实例的信息。要检查某一特定实例的信息，取消选择全部实例，再从列表中选择要显示的实例。

使用地图面板中的选项更改地图显示：
    \begin{itemize}
      \item Value：设置要显示的地图类型
      \item Bins：更改用于显示地图的唯一 bin 数量
      \item From: 与 To:：设置功耗密度地图中显示的密度值范围
      \item Text：在布局中显示计算值；若值不可见，放大以查看
    \end{itemize}
\end{enumerate}

\subsection{示例}
\subsubsection*{Examples}

下列示例说明如何在 GUI 中显示各类地图：
\begin{itemize}
  \item 图~\ref{fig:eff-vd-map}：显示轨实例有效压降地图并检查热点。
  \item 图~\ref{fig:switch-map}：显示 block 中开关单元的轨地图并检查所选开关单元的电压信息。
  \item 图~\ref{fig:power-dens}：显示 block 的功耗密度地图。
\end{itemize}

\begin{figure}[htbp]
\centering
\fbox{\parbox{0.85\textwidth}{\centering\vspace{1.5cm}
\textit{（原书 Figure 212：Displaying a Rail Instance Effective Voltage Drop Map）}\\[0.3em]
\small 选择 Rail Instance Effective Voltage Drop；悬停红色热点查看实例压降；Zoom to Critical；直方图与地图显示控制选项。请参阅原版 PDF 插图。
\vspace{1.5cm}}}
\caption{显示轨实例有效压降地图}
\label{fig:eff-vd-map}
\end{figure}

\begin{figure}[htbp]
\centering
\fbox{\parbox{0.85\textwidth}{\centering\vspace{1.5cm}
\textit{（原书 Figure 213：Displaying Maps for Switch Cells）}\\[0.3em]
\small 选择 Rail Instance Switch Cell；选择地图类型与实例；悬停查看电压信息。请参阅原版 PDF 插图。
\vspace{1.5cm}}}
\caption{显示开关单元地图}
\label{fig:switch-map}
\end{figure}

\begin{figure}[htbp]
\centering
\fbox{\parbox{0.85\textwidth}{\centering\vspace{1.5cm}
\textit{（原书 Figure 214：Displaying a Power Density Map）}\\[0.3em]
\small 选择 Rail Power，再选择 Density 显示功耗密度值。请参阅原版 PDF 插图。
\vspace{1.5cm}}}
\caption{显示功耗密度地图}
\label{fig:power-dens}
\end{figure}

\begin{seeAlsoBox}
指定 RedHawk 与 RedHawk-SC 工作目录
\end{seeAlsoBox}

% ============================================================
\section{在 GUI 中显示 ECO 形状}
\subsection*{Displaying ECO Shapes in the GUI}

当使用 RedHawk \texttt{mesh add} 命令在顶层 block 中为当前子 block 添加虚拟 PG mesh 时，默认这些虚拟 PG mesh 不会作为布局对象显示在 GUI 中。因此，这些 GUI 形状无法用布局对象集合命令（例如 \cmd{get\_shapes}）查询，保存 block 时也不会保存到设计库。

要将 ECO 形状与轨结果一起保存并在 GUI 中显示，启用 \appopt{rail.display\_eco\_shapes}：
\begin{lstlisting}
set_app_options -name rail.display_eco_shapes -value true
\end{lstlisting}

\begin{noteBox}
\texttt{mesh add} 命令仅在 RedHawk 中可用，不在 RedHawk-SC 中可用。因此，若同时启用 \appopt{rail.enable\_redhawk\_sc} 与 \appopt{rail.display\_eco\_shapes}，工具会发出错误消息。

要在 GUI 中显示 ECO 形状，必须先将 \appopt{rail.allow\_redhawk\_license\_checkout} 设为 \texttt{true} 以启用 RedHawk 签核许可证密钥，否则发出错误消息。
\end{noteBox}

要添加虚拟 PG mesh 并在 GUI 中显示所创建的 ECO 形状：
\begin{enumerate}
  \item 准备包含 \texttt{mesh add} 命令的脚本文件（例如 \texttt{mesh.tcl}）。
  \item 使用 \texttt{mesh.tcl} 脚本文件运行设计设置：
\begin{lstlisting}
fc_shell> set_rail_command_options -script_file mesh.tcl \
   -command setup_design -order after_the_command
\end{lstlisting}
工具生成 RedHawk 脚本文件，并在设计设置之后 source \texttt{mesh.tcl}。
  \item 使用下列命令执行轨分析：
\begin{lstlisting}
fc_shell> analyze_rail -nets -voltage_drop
\end{lstlisting}
  \item 运行 \cmd{open\_rail\_result} 加载设计数据与分析结果。
  \item 运行 \cmd{gui\_start} 打开 Fusion Compiler GUI。在 GUI 布局窗口中选择 View $>$ Map 并选择要显示的地图。
\end{enumerate}

此时可在 GUI 中显示并查询由 RedHawk \texttt{mesh add} 创建的 ECO 形状。

\subsection{脚本示例}
\subsubsection*{Script Examples}

使用 \cmd{set\_rail\_command\_options} 在 GUI 中显示 ECO 形状：
\begin{lstlisting}
## Open design ##
open_block
link_block
## Specify taps ##
create_taps
## Specify RedHawk Fusion input files or variables ##
set_app_options -name rail.enable_redhawk -value 1
set_app_options -name rail.redhawk_path -value
set_app_options -name rail.disable_timestamps -value true
set_app_options -name rail.display_eco_shapes -value true
set_rail_command_options -script_file mesh.tcl -command \
   setup_design ...
## Analyze##
analyze_rail -voltage_drop ... -nets {VDD VSS}
## Check GUI ##
open_rail_result
\end{lstlisting}

使用 RedHawk 脚本在 GUI 中显示 ECO 形状：
\begin{lstlisting}
## Open design ##
open_block
link_block
## Specify taps ##
create_taps
## Specify RedHawk Fusion input files or variables ##
set_app_options -name rail.enable_redhawk -value 1
set_app_options -name rail.redhawk_path -value
set_app_options -name rail.disable_timestamps -value true
set_app_options -name rail.display_eco_shapes -value true
## Analyze ##
analyze_rail -voltage_drop ... -nets {VDD VSS} -script_only
## Modify the script.tcl containing mesh add command ##
analyze_rail -redhawk_script_file
## Check GUI ##
open_rail_result
\end{lstlisting}

% ============================================================
\section{电压热点分析}
\subsection*{Voltage Hotspot Analysis}

轨分析完成后，分析报告中可能报告数百到数千个压降违例，从而难以识别违例的根本原因。

使用电压热点分析能力，通过将整个芯片划分为小网格框为电源或地 net 生成热点，并报告所生成网格框的轨相关信息。此外，可使用 \cmd{get\_instance\_result} 与 \cmd{get\_geometry\_result} 仅查询特定分析类型的基于单元或基于几何的轨结果。

热点分析提供下列功能：
\begin{itemize}
  \item 识别高压降 aggressor 的根本原因，例如 aggressor 自身的大电流，或因时序窗重叠导致同时翻转的邻近单元的电流。
  \item 根据报告确定用于修复电压违例的设计技术，例如选择候选参考单元进行尺寸替换、移动或重定位单元以降低压降，或应用 PG 增强以解决压降问题。
\end{itemize}

本节包含下列主题：
\begin{itemize}
  \item 生成热点（Generating Hotspots）
  \item 报告热点（Reporting Hotspots）
  \item 移除热点（Removing Hotspots）
  \item 电压热点分析示例（Voltage Hotspot Analysis Examples）
\end{itemize}

要仅写出位于热点区域内的单元实例或几何的轨结果，使用 \cmd{get\_instance\_result} 与 \cmd{get\_geometry\_result}。详见「生成基于实例的分析报告」与「生成基于几何的分析报告」。

\subsection{生成热点}
\subsubsection*{Generating Hotspots}

要为电源或地 net 生成压降热点，使用 \cmd{generate\_hot\_spots}。该命令将整个芯片区域按行与列划分为多个网格（例如 100$\times$100）。网格框按每个网格框中的最大有效压降值排序，并用从 0 开始的整数索引。

要基于标准单元 site row 创建行，并用最低金属层上的垂直 PG strap 创建列，使用 \opt{-site\_row}。或者，用 \opt{-row} 与 \opt{-column} 指定划分整个区域的行数与列数。

命令为压降值大于理想压降指定百分比的网格创建错误单元。默认为 0.1，即对压降值大于理想电源电压 10\% 的网格创建错误单元。错误单元类型为有效压降。

要检查生成的错误单元内容，在错误浏览器中打开错误单元。要检查网格框内容，使用 \cmd{report\_hot\_spots}。要从内存中移除生成的电压热点数据，使用 \cmd{remove\_hot\_spots}。

\begin{tabularx}{\textwidth}{@{}l X@{}}
\toprule
\textbf{选项} & \textbf{说明} \\
\midrule
\opt{-net} & 指定生成热点的 net 名。未指定时使用全部电源 net。 \\
\opt{-percentage} & 指定用于生成热点错误单元的百分比。默认为 0.1。 \\
\opt{-site\_row} \textit{integer} & 默认使用标准 site row 高度将整个区域划分为多行，用最低 PG net 垂直 strap 划分为多列。要使用 site row 高度的倍数创建网格行，指定带整数的 \opt{-site\_row}。例如设为 2 时，每两个 site row 构造一行。 \\
\opt{-row} & 指定划分整个区域的行数。 \\
\opt{-column} & 指定划分整个区域的列数。 \\
\bottomrule
\end{tabularx}

\begin{seeAlsoBox}
报告热点；移除热点；电压热点分析示例
\end{seeAlsoBox}

\subsection{报告热点}
\subsubsection*{Reporting Hotspots}

使用 \cmd{generate\_hot\_spots} 在 block 上生成热点后，运行 \cmd{report\_hot\_spots}，报告热点网格框中带电压相关单元属性的单元。电压相关单元属性包括：\texttt{current}、\texttt{overlap\_current}、\texttt{effective\_resistance}、\texttt{timing\_window}、\texttt{slack}、\texttt{output\_load} 与 \texttt{static\_power}。

使用这些电压相关单元属性识别电压违例的根本原因。例如，可以：
\begin{itemize}
  \item 使用电流与输出负载电容信息，判断单元的高峰值是否由负载过大引起。若是，拆分输出可能是降低高压降的有效方法。
  \item 检查并比较当前热点与邻近网格的时序窗，判断单元是否可移到其它网格（目标）。这样做可降低原网格的压降，而不增加目标网格的压降。
  \item 比较区域内全部单元的 slew 与负载电容，判断高峰值电流是否由陡峭 slew 或大负载电容引起。
  \item 查找时序窗重叠的单元，识别动态轨分析期间可能同时翻转的单元。
\end{itemize}

\cmd{report\_hot\_spots} 的常用选项：

\begin{tabularx}{\textwidth}{@{}l X@{}}
\toprule
\textbf{选项} & \textbf{说明} \\
\midrule
\opt{-index} & 按索引列出热点网格。 \\
\opt{-summary} & 报告全部网格的摘要，例如创建的网格数、block 中的实例数，以及压降值高于理想压降指定百分比的网格数。 \\
\opt{-object} & 指定要报告的对象名；列出包含指定单元实例对象的热点网格。 \\
\opt{-verbose} & 打印详细报告。 \\
\bottomrule
\end{tabularx}

\subsubsection{按索引列出电压网格框}
\subsubsection*{Listing Voltage Grid Boxes by Index}

默认情况下，工具向上下各搜索三行、向左右各搜索三列，以确定当前网格的邻近网格。使用 \opt{-index} 指定用于报告电压热点网格的索引。

索引号最小的网格框具有最高的有效压降值。从而可确定热点区域的大小。例如，若周围网格索引号都较大，而仅少数中心网格索引号小得多，则假定该压降热点为岛状（island）。

\subsubsection{累积电流}
\subsubsection*{Accumulating Current}

对于网格框中的实例，通过累积其邻近单元上基于时序窗的电流来估计重叠电流（overlap current）。

命令按时序窗从小到大扫描，并随时间流逝累积电流。利用该信息可判断有多少单元同时翻转，以及为降低压降偏移时序窗时应施加多少时间偏移。

假定设计有三个实例及其电流如下。如下列报告所示，点 2.0 处的累积电流来自 \texttt{inst2} 与 \texttt{inst3}，不包含 \texttt{inst1} 的电流：
\begin{lstlisting}
inst1 [0.0 2.0] current 0.11
inst2 [1.0 2.5] current 0.22
inst3 [2.0 3.0] current 0.33
\end{lstlisting}

\begin{seeAlsoBox}
生成热点；移除热点；电压热点分析示例
\end{seeAlsoBox}

\subsection{移除热点}
\subsubsection*{Removing Hotspots}

工具将生成的热点数据保存在内存中。若要在另一次运行中为其它 net 生成热点，需先用 \cmd{remove\_hot\_spots} 从内存中移除先前生成的热点数据，再继续另一次热点分析。

例如，若已对 VDD net 运行 \cmd{generate\_hot\_spots}，则必须先移除 VDD net 的热点数据，再为其它 net 生成热点。

\begin{seeAlsoBox}
生成热点；报告热点；电压热点分析示例
\end{seeAlsoBox}

\subsection{电压热点分析示例}
\subsubsection*{Voltage Hotspot Analysis Examples}

本主题提供如何使用热点分析能力识别电压违例根本原因的示例。

时序窗重叠的单元可能同时翻转，从而导致高压降。图~\ref{fig:hotspot-tw1} 说明如何识别单元是否具有重叠时序窗。

\begin{figure}[htbp]
\centering
\fbox{\parbox{0.85\textwidth}{\centering\vspace{1.5cm}
\textit{（原书 Figure 215：Checking High Voltage Drop Caused by Overlapping Timing Window (I)）}\\[0.3em]
\small 请参阅原版 PDF 插图。
\vspace{1.5cm}}}
\caption{检查由重叠时序窗引起的高压降（I）}
\label{fig:hotspot-tw1}
\end{figure}

接着，运行带 \opt{-index} 的 \cmd{report\_hot\_spots}，查找哪些网格具有导致大电流的重叠时序窗。如图~\ref{fig:hotspot-tw2} 所示，当 \opt{-index} 设为 5 时，网格 5 报告的总电流较小；当 \opt{-index} 设为 2 时，网格 2 报告的总电流较大。此时，重叠时序窗导致网格 2 中的有效压降较高。

\begin{figure}[htbp]
\centering
\fbox{\parbox{0.85\textwidth}{\centering\vspace{1.5cm}
\textit{（原书 Figure 216：Checking High Voltage Drop Caused by Overlapping Timing Window (II)）}\\[0.3em]
\small 请参阅原版 PDF 插图。
\vspace{1.5cm}}}
\caption{检查由重叠时序窗引起的高压降（II）}
\label{fig:hotspot-tw2}
\end{figure}

\begin{seeAlsoBox}
生成热点；报告热点；移除热点
\end{seeAlsoBox}

% ============================================================
\section{查询属性}
\subsection*{Querying Attributes}

轨分析或缺失 via 检查完成后，使用 \cmd{get\_attribute} 查询存储在 RedHawk 工作目录中的结果。可查询表~\ref{tab:rail-attrs} 所示的、专用于轨分析结果的属性。

\begin{table}[htbp]
\centering
\caption{轨结果对象属性（原书 Table 71）}
\label{tab:rail-attrs}
\small
\begin{tabularx}{\textwidth}{@{}l X@{}}
\toprule
\textbf{对象} & \textbf{属性} \\
\midrule
Cell & \texttt{static\_power} \\
\midrule
Pin &
\texttt{static\_current}；
\texttt{peak\_current}；
\texttt{static\_power}；
\texttt{switching\_power}；
\texttt{leakage\_power}；
\texttt{internal\_power}；
\texttt{min\_path\_resistance}；
\texttt{voltage\_drop}；
\texttt{average\_effective\_voltage\_drop\_in\_tw}；
\texttt{max\_effective\_voltage\_drop}（$=$ \texttt{effective\_voltage\_drop}）；
\texttt{max\_effective\_voltage\_drop\_in\_tw}；
\texttt{min\_effective\_voltage\_drop\_in\_tw}；
\texttt{effective\_resistance} \\
\bottomrule
\end{tabularx}
\end{table}

下例检索 \texttt{cellA/VDD} pin 的 \texttt{effective\_voltage\_drop} 属性：
\begin{lstlisting}
fc_shell> get_attribute [get_pins cellA/VDD] effective_voltage_drop
-0.02812498807907104
\end{lstlisting}

下例检索 \texttt{cellA} 单元的 \texttt{static\_power} 属性：
\begin{lstlisting}
fc_shell> get_attribute [get_cells cellA] static_power
2.014179968645724e-05
\end{lstlisting}

% 原书第 12 章至此结束（约第 947 页）。
